\documentclass{situationist-preprint}

\title{The Two Drivers}
\msid{ETH-13-b}
\msfamily{ETH-13}
\msdiscipline{Normative and Applied Ethics}
\mstype{Companion manuscript}
\mskeywords{moral luck; culpability; scope of responsibility; reparation; agent-regret}

\begin{document}
\maketitle

\begin{abstract}
ETH-13-a resolves the problem of moral luck by distinguishing the degree of an agent's culpability from the scope of her responsibility. We work through the case of the two careless drivers in detail. Culpability fixes the blame each deserves and the self-reproach fitting for each; scope fixes what each is answerable for, and so the apologies, reparations and agent-regret that each owes or may fittingly feel. Adding a third, reckless but lucky, driver shows that the two dimensions vary independently, so that neither can be reduced to the other.
\end{abstract}

\section{Two dimensions, two sets of responses}

Two drivers drive equally carelessly; a child runs in front of the first and is killed. By Proposition~3.2 of ETH-13-a, the drivers are equally culpable and the first is responsible for more. Culpability is fixed by quality of will and evidence (Definition~2.3 and Principle~P1 of ETH-13-a); scope extends to the outcomes in which a controlled act issues through the realization of a risk it culpably created (Principle~P3 of ETH-13-a).

The distinction sorts our responses into two groups. \emph{Responses to culpability} are blame and self-reproach. Their intensity is fixed by the quality of will an act expresses, so it is the same for the two drivers. A spectator who knows both cases and blames the unlucky driver more is making a mistake, the mistake that Principle~P1 of ETH-13-a corrects. \emph{Responses to scope} are those directed at what an agent is answerable for: acknowledgment, apology, reparation and agent-regret in the sense of \citet{Williams1976}. These differ. The unlucky driver owes the child's parents an apology and, where possible, reparation, because the death is part of what he did; the lucky driver owes these to no one, since there is no one whose loss falls within the scope of his responsibility. His fitting response is self-reproach for his carelessness and resolve to drive better, which the unlucky driver also owes.

The legal distinction between careless driving and causing death by careless driving tracks scope. It is often criticized as punishing luck \citep{EnochMarmor2007}. On the present account, the additional penalty is justified only insofar as it expresses answerability for the outcome rather than greater culpability. Whether punishment should track scope at all is a question for the theory of punishment (ETH-26); what the distinction shows is that the law is responding to a real moral difference.

\section{A third driver}

Add a third driver, who drives recklessly, at high speed through a residential street, and reaches home without incident. His quality of will is worse than that of either careless driver, so by Principle~P1 of ETH-13-a he is more culpable than both. The scope of his responsibility includes no death. He therefore exceeds the unlucky driver in culpability and falls below him in scope.

\begin{proposition}\label{p:indep}
Culpability and scope vary independently: there are agents $a$ and $b$ such that $a$ is more culpable than $b$ and $b$ is responsible for more than $a$.
\end{proposition}

\begin{proof}
Take $a$ to be the reckless driver and $b$ the unlucky careless driver. The comparisons were established above.
\end{proof}

The proposition rules out two reductions. One would measure responsibility by culpability alone, as the opponents of moral luck propose; it cannot say why the unlucky driver, and not the reckless one, owes the parents reparation. The other would measure blameworthiness by outcomes; it cannot say why the reckless driver deserves more blame than the unlucky one. Only the two-dimensional account captures both judgements.

\begin{objection}
Agent-regret is fitting for a driver who kills a child although he drove with all due care. Principle~P3 of ETH-13-a does not place the death within his scope, since he created no culpable risk; so scope cannot be what agent-regret tracks.
\end{objection}

\begin{reply}
ETH-13-a leaves open how far scope extends beyond Principle~P3, which states a sufficient condition. The blameless driver's regret is directed at his causal part in the death and may be fitting on that ground; but he owes no apology for wrongdoing, and any duty of reparation he has is weaker than the careless driver's. The difference between the two drivers is a difference that Principle~P3 marks, and it is enough for the account that the careless driver's agent-regret is fitting on the ground that the death lies within his scope.
\end{reply}

\section{Conclusion}

The two drivers deserve the same blame and owe different things. The reckless third driver deserves more blame and owes less. Moral luck operates on scope, and responsibility, whose condition is control over one's acts, remains in place.

\begin{thebibliography}{9}\small
\bibitem[Enoch and Marmor(2007)]{EnochMarmor2007} Enoch, D. and Marmor, A. (2007). The case against moral luck. \emph{Law and Philosophy}, 26(4), 405--436.
\bibitem[Williams(1976)]{Williams1976} Williams, B. (1976). Moral luck. \emph{Proceedings of the Aristotelian Society, Supplementary Volumes}, 50, 115--135.
\end{thebibliography}

\end{document}
